\chapter{Análise Pós-Laboratorial}

Nesta seção, realiza-se o tratamento matemático e a modelagem do filtro passa-baixa ativo de 4ª ordem a partir dos valores reais dos componentes passivos aferidos em bancada. O objetivo é reconstruir a função de transferência real do sistema, determinar as novas características de frequência de corte e ganho, e avaliar criticamente as discrepâncias observadas em relação aos dados nominais de projeto.

\section{Recálculo da Função de Transferência Real}

A função de transferência total do filtro bilinear de 4ª ordem é dada pelo produto das funções de transferência de cada um dos dois estágios Sallen-Key que compõem a estrutura em cascata, multiplicada pelo ganho de banda passante experimental ($K_{med}$):
\begin{equation}
	H_{real}(s) = K_{med} \cdot H_1(s) \cdot H_2(s) = K_{med} \cdot \left( \frac{1}{a_1 s^2 + b_1 s + 1} \right) \cdot \left( \frac{1}{a_2 s^2 + b_2 s + 1} \right)
\end{equation}

Os coeficientes polinomiais dependem diretamente do arranjo de componentes medidos a frio no laboratório. Utilizando as equações analíticas deduzidas na etapa teórica para componentes desiguais, os parâmetros para o \textbf{Estágio 1} são calculados como:
\begin{align}
	a_1 &= R_{1\_med} \cdot R_{2\_med} \cdot C_{1\_med} \cdot C_{2\_med} \nonumber \\
	a_1 &= 9998 \cdot 10056 \cdot (98,8 \times 10^{-9}) \cdot (97,7 \times 10^{-9}) \approx 9,7055 \times 10^{-7} \text{ s}^2 \\
	b_1 &= C_{2\_med} \cdot (R_{1\_med} + R_{2\_med}) \nonumber \\
	b_1 &= (97,7 \times 10^{-9}) \cdot (9998 + 10056) \approx 1,9593 \times 10^{-3} \text{ s}
\end{align}

Analogamente, para o \textbf{Estágio 2}, os coeficientes são determinados por:
\begin{align}
	a_2 &= R_{3\_med} \cdot R_{4\_med} \cdot C_{3\_med} \cdot C_{4\_med} \nonumber \\
	a_2 &= 9908 \cdot 10059 \cdot (102 \times 10^{-9}) \cdot (95 \times 10^{-9}) \approx 9,6571 \times 10^{-7} \text{ s}^2 \\
	b_2 &= C_{4\_med} \cdot (R_{3\_med} + R_{4\_med}) \nonumber \\
	b_2 &= (95 \times 10^{-9}) \cdot (9908 + 10059) \approx 1,8969 \times 10^{-3} \text{ s}
\end{align}

Substituindo os coeficientes numéricos obtidos, a função de transferência real do filtro no domínio de Laplace passa a ser expressa por:
\begin{equation}
	H_{real}(s) = \frac{5,6}{(9,7055 \times 10^{-7} s^2 + 1,9593 \times 10^{-3} s + 1)(9,6571 \times 10^{-7} s^2 + 1,8969 \times 10^{-3} s + 1)}
\end{equation}

\section{Análise Comparativa: Frequência de Corte e Ganho}

A partir do polinômio do denominador da função de transferência ajustada, pode-se extrair a nova frequência de corte teórica corrigida do filtro. Avaliando o ponto de atenuação de $-3\text{ dB}$ ($|H_{real}(j\omega)| = \frac{K_{med}}{\sqrt{2}}$) via \textit{software} algébrico \cite{maxima}, obtém-se uma frequência de corte teórica recalculada de aproximadamente $51,2\text{ Hz}$.

Ao comparar este valor com os dados nominais de projeto e com a medição prática realizada no laboratório, estabelece-se a análise de erros percentuais explicitada na Tabela \ref{tab:comparativo_erros}.

\begin{table}[htbp]
	\centering
	\caption{Comparativo global de parâmetros: Nominal, Recalculado e Medido.}
	\label{tab:comparativo_erros}
	\begin{tabular}{lcccc}
		\toprule
		\textbf{Parâmetro} & \textbf{Nominal} & \textbf{Recalculado} & \textbf{Medido} & \textbf{Erro} \\
		\midrule
		Freq. de Corte ($f_c$) & $50,0\text{ Hz}$ & $51,2\text{ Hz}$ & $46,0\text{ Hz}$ & $10,15\%$ \\
		Ganho de Banda ($K$) & $4,0\text{ V/V}$ & $4,0\text{ V/V}$ & $5,6\text{ V/V}$ & $40,00\%$ \\
		\bottomrule
	\end{tabular}
\end{table}

O erro associado à frequência de corte ($10,15\%$) encontra-se dentro do limite esperado para circuitos integrados com componentes discretos de tolerância comercial convencional (tipicamente $\pm 5\%$ para resistores e $\pm 10\%$ para capacitores). A variação cumulativa das capacitâncias reais (com desvios para menos no primeiro estágio e para mais no segundo) deslocou ligeiramente os polos originais da aproximação ideal de Butterworth, resultando na redução da frequência observada em bancada para $46\text{ Hz}$.

\section{Discussão sobre as Discrepâncias de Ganho e Não-Idealidades}

O desvio mais acentuado verificado na prática ocorreu no ganho contínuo da banda de passagem, que saltou de um valor projetado de $4,0$ para uma leitura nominal de $5,6$ em bancada. Esse fenômeno de sobreganho em filtros ativos Sallen-Key pode ser fundamentado por fatores físicos e instrumentais específicos:

\begin{enumerate}
	\item \textbf{Tolerância do Circuito de Realimentação:} O ganho de um estágio ativo Sallen-Key não unitário é ditado pela razão de resistores conectados à malha inversora do amplificador operacional \cite{nilsson}. Pequenas variações nas resistências internas de ganho ou o uso de resistores comerciais com desvios nominais cumulativos alteram diretamente o fator $K$, elevando a resposta de amplitude para sinais de baixa frequência.
	\item \textbf{Efeito de Carga e Cascateamento:} O modelo teórico assume isolamento perfeito entre o Estágio 1 e o Estágio 2. Na prática, a impedância de entrada do segundo estágio, combinada com a impedância de saída não nula em malha aberta do primeiro circuito integrado, causa um acoplamento dinâmico que distorce o fator de amortecimento global, elevando o pico de magnitude próximo à transição da banda de passagem.
	\item \textbf{Incerteza de Medição Instrumental:} Conforme observado na tabela de varredura, a amplitude da tensão de entrada ($V_i$) oscilou entre $4,8\text{ V}$ e $5,2\text{ V}$ devido à regulação de carga do gerador de sinais. Como o ganho prático é calculado de forma indireta pela razão $V_o/V_i$, variações instantâneas na amplitude de excitação inserem erros sistemáticos na quantificação precisa do ganho RMS real pelo osciloscópio.
\end{enumerate}

Em suma, as curvas e os pontos experimentais demonstram que, apesar das não-idealidades paramétricas terem elevado o ganho e deslocado sutilmente a frequência de corte, o circuito cumpriu com êxito a sua função principal de filtragem, apresentando o perfil de atenuação abrupto de quarta ordem esperado para frequências superiores a $50\text{ Hz}$.